\documentclass[11pt]{article}
\usepackage{mathblatt}
\usepackage{multicol}


\begin{document}
\pfheftstil
\pfheftkopf{Wahrscheinlichkeit \textperiodcentered{} Lösungen}{}
\begin{pfloesung}{}
\lz{1.}{AA, AB, BA, BB}{}{}
\end{pfloesung}
\begin{pfloesung}{}
\lz{2.}{$\tfrac{1}{20}$ = 5 \%; $\tfrac{1}{5}$ = 20 \%; $\tfrac{1}{3}$ $\approx$ 33,3 \%; $\tfrac{3}{20}$ = 15 \%}{}{}
\end{pfloesung}
\begin{pfloesung}{}
\lz{3.}{$\tfrac{7}{8}$; $\tfrac{1}{2}$; 0,3}{}{}
\end{pfloesung}
\begin{pfloesung}{}
\lz{4.}{$\tfrac{1}{16}$; $\tfrac{1}{4}$; $\tfrac{5}{16}$}{}{}
\end{pfloesung}
\begin{pfloesung}{}
\lz{5.}{15; 2}{}{}
\end{pfloesung}
\begin{pfloesung}{}
\lz{6.}{4 Paare: , , ,}{grau; gelb; grau; weiß; grün; gelb; grün; weiß}{}
\end{pfloesung}
\begin{pfloesung}{}
\lz{7.}{$6$ Kombinationen}{$3 \cdot 2$ $\Rightarrow$ 6}{}
\end{pfloesung}
\begin{pfloesung}{}
\lz{8.}{123, 132, 213, 231, 312, 321}{}{}
\end{pfloesung}
\begin{pfloesung}{}
\lz{9.}{(1,5), (2,4), (3,3), (4,2), (5,1)}{}{}
\end{pfloesung}
\begin{pfloesung}{}
\lz{10.}{6 Paare: , , , , ,}{1; 4; 2; 4; 3; 4; 4; 4; 5; 4; 4}{}
\end{pfloesung}
\begin{pfloesung}{}
\lz{11.}{3 · 2 · 4 = 24 Möglichkeiten}{}{}
\end{pfloesung}
\begin{pfloesung}{}
\lz{12.}{5 · 4 · 3 · 2 · 1 = 120 > 100 – stimmt}{}{}
\end{pfloesung}
\begin{pfloesung}{}
\lz{13.}{6 Möglichkeiten (4 · 3 : 2)}{}{}
\end{pfloesung}
\begin{pfloesung}{}
\lz{14.}{4 · 3 · 2 · 1 = 24}{}{}
\end{pfloesung}
\begin{pfloesung}{}
\lz{15.}{wr, wb, wg, rb, rg, bg (6 Paare)}{}{}
\end{pfloesung}
\begin{pfloesung}{}
\lz{16.}{9 Zahlen: 44, 45, 46, 54, 55, 56, 64, 65, 66}{}{}
\end{pfloesung}
\begin{pfloesung}{}
\lz{17.}{4$^{6}$ = 4 096}{}{}
\end{pfloesung}
\begin{pfloesung}{}
\lz{18.}{11, 12, 13, 21, 22, 23, 31, 32, 33}{3 · 3 $\Rightarrow$ 9 Zahlen}{}
\end{pfloesung}
\begin{pfloesung}{}
\lz{19.}{4 + 3 + 2 + 1 = 10}{}{}
\end{pfloesung}
\begin{pfloesung}{}
\lz{20.}{$\tfrac{1}{3}$}{6 Zahlen: 147, 174, 417, 471, 714, 741; gerade sind: $174$ und $714,$ Anteil $\frac{2}{6}$}{}
\end{pfloesung}
\begin{pfloesung}{}
\lz{21.}{9 Ergebnisse: , , , , , , , ,}{5; 5; 5; 7; 5; 7; 5; 7; 7; 7; 5; 7}{}
\end{pfloesung}
\begin{pfloesung}{}
\lz{22.}{3 Ergebnisse: , ,}{Kopf; Kopf; Kopf; Zahl; Zahl; Kopf}{}
\end{pfloesung}
\begin{pfloesung}{}
\lz{23.}{4 (ZZ, ZW, WZ, WW)}{}{}
\end{pfloesung}
\begin{pfloesung}{}
\lz{24.}{(1 | 2), (2 | 2), (3 | 2), (4 | 2), (5 | 2), (6 | 2); $\tfrac{1}{6}$ $\approx$ 16,7 \%}{6 von 36 Augenpaaren}{}
\end{pfloesung}
\begin{pfloesung}{}
\lz{25.}{$\tfrac{1}{2}$}{}{}
\end{pfloesung}
\begin{pfloesung}{}
\lz{26.}{möglich: 35 Lose , günstig: 5 Gewinne}{30 Nieten und 5 Gewinne}{}
\end{pfloesung}
\begin{pfloesung}{}
\lz{27.}{eine Stufe}{}{}
\end{pfloesung}
\begin{pfloesung}{}
\lz{28.}{Nein: Beide Augenzahlen haben die Wahrscheinlichkeit $\tfrac{1}{6}$}{}{}
\end{pfloesung}
\begin{pfloesung}{}
\lz{29.}{$\tfrac{1}{2}$}{$\tfrac{3}{6}$ $\Rightarrow$ 0{,}5}{}
\end{pfloesung}
\begin{pfloesung}{}
\lz{30.}{$\tfrac{1}{5}$ = 20 \%}{}{}
\end{pfloesung}
\begin{pfloesung}{}
\lz{31.}{$\tfrac{2}{3}$}{}{}
\end{pfloesung}
\begin{pfloesung}{}
\lz{32.}{$\tfrac{1}{20}$ = 5 \%}{}{}
\end{pfloesung}
\begin{pfloesung}{}
\lz{33.}{$\tfrac{3}{20}$ = 15 \%}{}{}
\end{pfloesung}
\begin{pfloesung}{}
\lz{34.}{$\tfrac{1}{5}$ = 20 \%}{1 günstiger von 5 Zetteln}{}
\end{pfloesung}
\begin{pfloesung}{}
\lz{35.}{mittlere Dose}{}{}
\end{pfloesung}
\begin{pfloesung}{}
\lz{36.}{rechter Topf (3 weiße, 3 graue Kugeln)}{}{}
\end{pfloesung}
\begin{pfloesung}{}
\lz{37.}{Nein: $\tfrac{3}{16}$, denn es sind nur noch 16 Muffins da, davon 3 mit Chili}{}{}
\end{pfloesung}
\begin{pfloesung}{}
\lz{38.}{$18\,\%$}{$\tfrac{9}{50}$ $\Rightarrow$ 0{,}18}{}
\end{pfloesung}
\begin{pfloesung}{}
\lz{39.}{$\tfrac{1}{10}$ = 10 \%}{1 von 10 Ziffern}{}
\end{pfloesung}
\begin{pfloesung}{}
\lz{40.}{Kopf: $\tfrac{1}{2}$ > $\tfrac{1}{6}$}{}{}
\end{pfloesung}
\begin{pfloesung}{}
\lz{41.}{$0{,}35$}{$\tfrac{35}{100}$ $\Rightarrow$ $\tfrac{7}{20}$}{}
\end{pfloesung}
\begin{pfloesung}{}
\lz{42.}{900 $-$ 101 + 1 = 800 Lose, 1 Hauptgewinn $\Rightarrow$ $\tfrac{1}{800}$}{}{}
\end{pfloesung}
\begin{pfloesung}{}
\lz{43.}{$\tfrac{3}{8}$; 37,5 \%}{3 von 8 Feldern}{}
\end{pfloesung}
\begin{pfloesung}{}
\lz{44.}{nein}{P = $\tfrac{2}{14}$ = $\tfrac{1}{7}$ $\approx$ 14,3 \%, nicht $\tfrac{2}{16}$; noch: 14 Pfannkuchen, 2 mit Senf}{}
\end{pfloesung}
\begin{pfloesung}{}
\lz{45.}{ja}{7 Trostpreis-Lose unter 80 Losen mit Endziffer 6 $\Rightarrow$ $\tfrac{7}{80}$; 80 Lose mit Endziffer 6; 8 mit Endung 26, ohne 326 = 7}{}
\end{pfloesung}
\begin{pfloesung}{}
\lz{46.}{8 Kekse, davon 3 mit Schoko}{}{}
\end{pfloesung}
\begin{pfloesung}{}
\lz{47.}{Gegenereignis: keine Sechs in beiden Würfen}{}{}
\end{pfloesung}
\begin{pfloesung}{}
\lz{48.}{mit Zurücklegen}{}{}
\end{pfloesung}
\begin{pfloesung}{}
\lz{49.}{2 Pfade: ,}{rot; blau; blau; rot}{}
\end{pfloesung}
\begin{pfloesung}{}
\lz{50.}{auf beiden Stufen rot $\tfrac{1}{3}$, blau $\tfrac{2}{3}$}{}{}
\end{pfloesung}
\begin{pfloesung}{}
\lz{51.}{$\tfrac{3}{28}$}{$\tfrac{3}{8} \cdot \tfrac{2}{7}$ $\Rightarrow$ $\tfrac{6}{56}$}{}
\end{pfloesung}
\begin{pfloesung}{}
\lz{52.}{zwei Stufen, je Ast 2: ½, 3: $\tfrac{2}{5}$, 5: 1/10, in der zweiten Stufe gleich}{}{}
\end{pfloesung}
\begin{pfloesung}{}
\lz{53.}{nach rot: rot $\tfrac{3}{6}$}{erste Stufe grün: $\tfrac{3}{7}$ $\Rightarrow$ 0{,}43}{}
\end{pfloesung}
\begin{pfloesung}{}
\lz{54.}{1. Woche: kein Joker 4/5, Joker $\tfrac{1}{5}$; 2. Woche an beiden Ästen: kein Joker 4/5, Joker $\tfrac{1}{5}$; P(Joker, Joker) = $\tfrac{1}{5}$ · $\tfrac{1}{5}$ = $\tfrac{1}{25}$}{mit Zurücklegen: jede Woche $\tfrac{1}{5}$}{}
\end{pfloesung}
\begin{pfloesung}{}
\lz{55.}{Baum mit $\tfrac{3}{10}$ und $\tfrac{7}{10}$ auf jeder Stufe; zweimal weiß: $\tfrac{7}{10}$ · $\tfrac{7}{10}$ = 49 \%}{}{}
\end{pfloesung}
\begin{pfloesung}{}
\lz{56.}{z. B. ein Glücksrad mit fünf gleich großen Feldern, zwei rot und drei blau, zweimal gedreht}{}{}
\end{pfloesung}
\begin{pfloesung}{}
\lz{57.}{1. Stufe A 1/2, B $\tfrac{3}{8}$; 2. Stufe je Ast A 1/2, B 3/8, C $\tfrac{1}{8}$}{}{}
\end{pfloesung}
\begin{pfloesung}{}
\lz{58.}{Stufe 1: T 0,35; H 0,65. Stufe 2: nach T: F 0,18, oF 0,82; nach H: F p/100, oF 1 $-$ p/100}{}{}
\end{pfloesung}
\begin{pfloesung}{}
\lz{59.}{$\tfrac{16}{20}$; $\tfrac{3}{18}$; Aussage 1 falsch, 2 falsch, 3 wahr}{fetter Pfad = dreimal rot; $\tfrac{4}{20}$ · $\tfrac{3}{19}$ · $\tfrac{2}{18}$ $\Rightarrow$ $\tfrac{1}{285}$ $\approx$ 0,35 \%}{}
\end{pfloesung}
\begin{pfloesung}{}
\lz{60.}{Äste 6: 1/6, keine 6: $\tfrac{5}{6}$; $\tfrac{2}{27}$ $\approx$ 7,4 \%}{3 · ($\tfrac{1}{6}$)² · $\tfrac{5}{6}$ $\Rightarrow$ $\tfrac{15}{216}$; ($\tfrac{1}{6}$)³ $\Rightarrow$ $\tfrac{1}{216}$}{}
\end{pfloesung}
\begin{pfloesung}{}
\lz{61.}{10/11, 1/10, 1/9, $\tfrac{8}{9}$; $\tfrac{3}{11}$ $\approx$ 0,27}{Unterschied: 1 $-$ $\tfrac{10}{11}$ · $\tfrac{9}{10}$ · $\tfrac{8}{9}$ $\Rightarrow$ 0{,}27; jeder Pfad mit R hat: $\tfrac{1}{11}$ $\Rightarrow$ 0{,}09}{}
\end{pfloesung}
\begin{pfloesung}{}
\lz{62.}{½ · ½ = $\tfrac{1}{4}$}{}{}
\end{pfloesung}
\begin{pfloesung}{}
\lz{63.}{$\tfrac{6}{10}$ · $\tfrac{1}{2}$ = $\tfrac{3}{10}$ = 30 \%}{}{}
\end{pfloesung}
\begin{pfloesung}{}
\lz{64.}{P = $\tfrac{3}{10}$ · $\tfrac{3}{10}$ = $\tfrac{9}{100}$}{}{}
\end{pfloesung}
\begin{pfloesung}{}
\lz{65.}{$\tfrac{9}{64}$}{$\tfrac{3}{8} \cdot \tfrac{3}{8}$ $\Rightarrow$ 0{,}14}{}
\end{pfloesung}
\begin{pfloesung}{}
\lz{66.}{$\tfrac{7}{16}$}{$\tfrac{3}{4} \cdot \tfrac{2}{4} + \tfrac{1}{4} \cdot \tfrac{1}{4}$ $\Rightarrow$ $\tfrac{6}{16} + \tfrac{1}{16}$}{}
\end{pfloesung}
\begin{pfloesung}{}
\lz{67.}{Nein: Zwei gleiche Augenzahlen haben $\tfrac{6}{36} = \tfrac{1}{6}$, eine 1 im ersten Wurf hat auch $\tfrac{1}{6}$}{}{}
\end{pfloesung}
\begin{pfloesung}{}
\lz{68.}{Nein. Der Gutschein wird zurückgelegt; für Jarne bleibt P(K) = $\tfrac{1}{10}$.}{}{}
\end{pfloesung}
\begin{pfloesung}{}
\lz{69.}{Ja. $\tfrac{5}{20}$ · $\tfrac{3}{20}$ = $\tfrac{3}{20}$ · $\tfrac{5}{20}$ = $\tfrac{15}{400}$.}{}{}
\end{pfloesung}
\begin{pfloesung}{}
\lz{70.}{P = $\tfrac{1}{3}$ · $\tfrac{1}{3}$ = $\tfrac{1}{9}$}{}{}
\end{pfloesung}
\begin{pfloesung}{}
\lz{71.}{$0{,}64$}{neuer Wert: $0{,}8 \cdot 0{,}8$ $\Rightarrow$ 0{,}64}{}
\end{pfloesung}
\begin{pfloesung}{}
\lz{72.}{P(33) = $\tfrac{1}{4}$ · $\tfrac{1}{4}$ = $\tfrac{1}{16}$; falsch, P(22) = $\tfrac{1}{4}$ · $\tfrac{2}{4}$ = $\tfrac{1}{8}$ $\neq$ $\tfrac{1}{16}$}{P(3 links) = $\tfrac{1}{4}$; P(2 rechts) = $\tfrac{2}{4}$}{}
\end{pfloesung}
\begin{pfloesung}{}
\lz{73.}{gleich: $\tfrac{36}{100}$ + $\tfrac{9}{100}$ + $\tfrac{1}{100}$ = $\tfrac{46}{100}$; verschieden: $\tfrac{54}{100}$; zwei verschiedene sind wahrscheinlicher.}{}{}
\end{pfloesung}
\begin{pfloesung}{}
\lz{74.}{Max hat nicht Recht; beide Wahrscheinlichkeiten $\tfrac{1}{6}$}{$\tfrac{6}{36}$ $\Rightarrow$ $\tfrac{1}{6}$; P(zuerst 6) = $\tfrac{1}{6}$}{}
\end{pfloesung}
\begin{pfloesung}{}
\lz{75.}{$\tfrac{5}{16}$ = 31,25 \%}{P(12) = $\tfrac{4}{16}$; P(21) = $\tfrac{1}{16}$}{}
\end{pfloesung}
\begin{pfloesung}{}
\lz{76.}{0,288² + 0,152² + 0,196² + 0,03² + 0,146² + 0,188² $\approx$ 0,202 $\approx$ 20 \%}{}{}
\end{pfloesung}
\begin{pfloesung}{}
\lz{77.}{$\tfrac{13}{32}$ $\approx$ 0,41}{$\tfrac{1}{4}$ + $\tfrac{9}{64}$ + $\tfrac{1}{64}$ $\Rightarrow$ $\tfrac{26}{64}$}{}
\end{pfloesung}
\begin{pfloesung}{}
\lz{78.}{$\tfrac{2}{10}$ · $\tfrac{1}{9}$ = $\tfrac{1}{45}$ $\approx$ 2,2 \%}{}{}
\end{pfloesung}
\begin{pfloesung}{}
\lz{79.}{P = $\tfrac{9}{18}$ = $\tfrac{1}{2}$}{}{}
\end{pfloesung}
\begin{pfloesung}{}
\lz{80.}{P = $\tfrac{11}{47}$}{}{}
\end{pfloesung}
\begin{pfloesung}{}
\lz{81.}{Recht: gleich viele Hauptgewinne, aber nur noch 460 Lose}{}{}
\end{pfloesung}
\begin{pfloesung}{}
\lz{82.}{Nach dem ersten Zug liegen nur noch 7 Kugeln in der Schachtel.}{}{}
\end{pfloesung}
\begin{pfloesung}{}
\lz{83.}{$\tfrac{2}{5}$}{$\tfrac{3}{5} \cdot \tfrac{2}{4} + \tfrac{2}{5} \cdot \tfrac{1}{4}$ $\Rightarrow$ $\tfrac{6}{20} + \tfrac{2}{20}$ = $\tfrac{8}{20}$}{}
\end{pfloesung}
\begin{pfloesung}{}
\lz{84.}{$\tfrac{1}{3}$}{$\tfrac{1}{9} + \tfrac{8}{9} \cdot \tfrac{1}{8} + \tfrac{8}{9} \cdot \tfrac{7}{8} \cdot \tfrac{1}{7}$ $\Rightarrow$ $\tfrac{3}{9}$}{}
\end{pfloesung}
\begin{pfloesung}{}
\lz{85.}{zehnmal so wahrscheinlich, nicht doppelt}{Drei rote: $\tfrac{5}{15} \cdot \tfrac{4}{14} \cdot \tfrac{3}{13}$ $\Rightarrow$ $\frac{60}{2\,730},$ drei grüne: $\frac{3}{15} \cdot \frac{2}{14} \cdot \frac{1}{13}$ = $\tfrac{6}{2\,730}$}{}
\end{pfloesung}
\begin{pfloesung}{}
\lz{86.}{$0{,}7$}{Unterschied: $1 - \tfrac{3}{5} \cdot \tfrac{2}{4}$ $\Rightarrow$ $1 - \tfrac{6}{20}$}{}
\end{pfloesung}
\begin{pfloesung}{}
\lz{87.}{beim 21. Zug: nach 20 A-Kugeln sind nur noch B-Kugeln in der Urne (P = 100 \%)}{}{}
\end{pfloesung}
\begin{pfloesung}{}
\lz{88.}{P(ss) = $\tfrac{4}{6}$ · $\tfrac{3}{5}$ = $\tfrac{12}{30}$ = $\tfrac{2}{5}$ = 40 \%}{}{}
\end{pfloesung}
\begin{pfloesung}{}
\lz{89.}{$\tfrac{1}{5}$}{$\tfrac{4}{5}$ · $\tfrac{1}{4}$ $\Rightarrow$ $\tfrac{4}{20}$}{}
\end{pfloesung}
\begin{pfloesung}{}
\lz{90.}{Ja. Vorher $\tfrac{4}{16}$ = 25 \%, jetzt $\tfrac{3}{15}$ = 20 \%.}{}{}
\end{pfloesung}
\begin{pfloesung}{}
\lz{91.}{a) $\tfrac{5}{20}$ · $\tfrac{4}{19}$ $\approx$ 5,26 \%; b) 1 $-$ $\tfrac{6}{20}$ · $\tfrac{5}{19}$ $\approx$ 92,1 \%; c) $\tfrac{11}{20}$ · $\tfrac{10}{19}$ $\approx$ 28,9 \%}{}{}
\end{pfloesung}
\begin{pfloesung}{}
\lz{92.}{$\tfrac{1}{57}$ $\approx$ 1,75 \% gegen $\tfrac{1}{1140}$ $\approx$ 0,09 \% $\Rightarrow$ 20-mal so hoch, Behauptung falsch}{$\tfrac{6}{20}$ · $\tfrac{5}{19}$ · $\tfrac{4}{18}$ $\Rightarrow$ $\tfrac{1}{57}$; $\tfrac{3}{20}$ · $\tfrac{2}{19}$ · $\tfrac{1}{18}$ $\Rightarrow$ $\tfrac{1}{1140}$}{}
\end{pfloesung}
\begin{pfloesung}{}
\lz{93.}{(S | M), (M | S), (S | S); $\tfrac{91}{120}$ $\approx$ 0,758; beide Pfannkuchen mit derselben Füllung}{$\tfrac{14}{16}$ · $\tfrac{13}{15}$ $\Rightarrow$ $\tfrac{182}{240}$}{}
\end{pfloesung}
\begin{pfloesung}{}
\lz{94.}{Ja}{„zwei verschiedene Zahlen“ ist das Gegenereignis; 100 \% $-$ 20 \% = 80 \%}{}
\end{pfloesung}
\begin{pfloesung}{}
\lz{95.}{1 $-$ $\tfrac{11}{30}$ = $\tfrac{19}{30}$ $\approx$ 63 \%}{}{}
\end{pfloesung}
\begin{pfloesung}{}
\lz{96.}{$\tfrac{13}{24}$}{}{}
\end{pfloesung}
\begin{pfloesung}{}
\lz{97.}{$\tfrac{2}{3}$}{}{}
\end{pfloesung}
\begin{pfloesung}{}
\lz{98.}{200 $-$ 70 $-$ 30 $-$ 10 = 90 Nieten}{}{}
\end{pfloesung}
\begin{pfloesung}{}
\lz{99.}{$\tfrac{7}{12}$}{Unterschied: $1 - \tfrac{3}{6} \cdot \tfrac{5}{6}$ $\Rightarrow$ $1 - \tfrac{15}{36}$ = $\tfrac{21}{36}$}{}
\end{pfloesung}
\begin{pfloesung}{}
\lz{100.}{1 $-$ ½ · ½ = $\tfrac{3}{4}$}{}{}
\end{pfloesung}
\begin{pfloesung}{}
\lz{101.}{1~ohne Zurücklegen, 2~mit Zurücklegen, 3~Pfadregel, 4~Gegenereignis, 5~ohne Zurücklegen, 6~mit Zurücklegen}{}{}
\end{pfloesung}
\begin{pfloesung}{}
\lz{102.}{z. B. Glücksrad mit 5 gleich großen Feldern, 1 Feld ist Gewinn: P = $\tfrac{1}{5}$ = 20 \%}{}{}
\end{pfloesung}
\begin{pfloesung}{}
\lz{103.}{$\tfrac{3}{8}$}{z. B. links: $5, 5, 5, 6$ und rechts $7, 7, 8, 8: \frac{3}{4} \cdot \frac{2}{4}$ $\Rightarrow$ $\tfrac{6}{16}$}{}
\end{pfloesung}
\begin{pfloesung}{}
\lz{104.}{2 rote Felder}{}{}
\end{pfloesung}
\begin{pfloesung}{}
\lz{105.}{z. B. 12 Murmeln: 3 blau, 2 grün, 6 gelb, 1 rot}{}{}
\end{pfloesung}
\begin{pfloesung}{}
\lz{106.}{r : (2r + g) = 0,2 $\Rightarrow$ g = 3r $\Rightarrow$ 3-mal so viele}{}{}
\end{pfloesung}
\begin{pfloesung}{}
\lz{107.}{z. B. links 1, 1, 2, 3 und rechts 3, 3, 1, 2: P(13) = $\tfrac{2}{4}$ · $\tfrac{2}{4}$ = $\tfrac{1}{4}$ = 25 \%}{$\tfrac{1}{4}$ $\Rightarrow$ $\tfrac{2}{4}$ · $\tfrac{2}{4}$ (oder 1 · $\tfrac{1}{4}$)}{}
\end{pfloesung}
\begin{pfloesung}{}
\lz{108.}{2 weiße Kugeln}{6 Kugeln insgesamt}{}
\end{pfloesung}
\begin{pfloesung}{\textbf{Prüfstein}}
\lz{a)}{$\tfrac{2}{6}$ = $\tfrac{1}{3}$}{2 günstige von 6 Flächen}{}
\lz{b)}{A ungerade $\tfrac{3}{6}$; B nach A gerade: ungerade $\tfrac{5}{6}$; B nach A ungerade: gerade 1/6, ungerade $\tfrac{5}{6}$; P(beide ungerade) = $\tfrac{3}{6}$ · $\tfrac{5}{6}$ = $\tfrac{15}{36}$ = $\tfrac{5}{12}$}{$\tfrac{1}{6}$ $\Rightarrow$ 0{,}17; $\tfrac{5}{6}$ $\Rightarrow$ 0{,}83}{}
\lz{c)}{$\tfrac{33}{36}$ = $\tfrac{11}{12}$}{$\tfrac{3}{6}$ · $\tfrac{1}{6}$ $\Rightarrow$ $\tfrac{3}{36}$}{}
\lz{d)}{z. B. B neu: 1, 1, 1, 3, 3, 3 (die 2 durch 1 ersetzt), P(1, 1) = $\tfrac{2}{6}$ · $\tfrac{3}{6}$ = $\tfrac{6}{36}$ = $\tfrac{1}{6}$}{Summe: 2 nur bei 1 + 1; 1) = $\tfrac{2}{6}$ = 1) = $\tfrac{3}{6}$}{}
\end{pfloesung}
\end{document}